\chapter{Discussion, Clinical Readiness Gap, Limitations, and Conclusion}
\label{ch:discussion_conclusion}

This thesis examined contactless neonatal respiratory-rate estimation as a chain of
measurement, signal processing, inference, and decision-making. The results do not support a
single claim that a high-accuracy estimator is ready to replace a bedside monitor. They support a
more bounded conclusion: in the evaluated cohort, a compact model with explicit signal and
temporal structure was more robust than the tested unconstrained regressors; simple uncertainty
scores and a fixed quality barrier exposed useful failure information; and independent phantom
measurements demonstrated controlled hardware transfer. At the same time, differences between
the available respiratory references, patient-level interval widths, and the absence of a
prospective clinical evaluation leave bedside utility unresolved.

The distinction matters because the output of a radar pipeline is not determined by its estimator
alone. Synchronization changes which samples are compared. A monitor trend and a rate computed
from individual respiratory triggers encode different timing and pause conventions. Filtering
sets the range of rates the system can represent. Finally, selective rejection and interval
calibration determine what a user sees when the signal becomes ambiguous. The following
discussion interprets the thesis across these connected stages and states the resulting claim
boundary explicitly.

\section{What the results establish}
\label{sec:ch7_synthesis}

\subsection{The value and limits of physical inductive bias}

The point-estimation results favor a structured, low-capacity model over the specific
high-capacity regressors evaluated here. On the 11-infant, 28-session \texttt{pool30\_sync}
cohort, the eight-parameter causal differentiable HMM R2e reduced equal-session RMSE relative to
the memoryless R0 baseline by 0.9687~bpm on the 1~Hz monitor reference (95\% interval
[$-1.1347,-0.8034$]; patient-level Holm-adjusted $p=0.0117$) and by 0.6635~bpm on the shared
\texttt{Trig -10 s} reference ($[-0.8833,-0.4464]$; Holm-adjusted $p=0.0215$). The paired
improvement appeared in 11/11 infants on the monitor reference and 10/11 on the lag-adjusted
trigger reference. This is evidence for R2e over R0 within the frozen evaluation protocol; it is
not evidence that the radar estimates a clinically adjudicated respiratory rate without error.

The comparison with the three-parameter forward-MMSE model R1a is more restrained. R2e's
one-hertz point estimate was 0.1418~bpm lower, but the patient-level comparison was not
significant after multiplicity correction ($p_{\mathrm{pat}}=0.0830$;
$p_{\mathrm{pat,Holm}}=0.6641$). The patient-level comparison on \texttt{Trig -10 s} was also
non-significant ($p_{\mathrm{pat}}=0.5771$). Thus the evidence supports a useful gain over the
memoryless baseline, while the cohort does not resolve whether eight parameters improve on the
much smaller temporal model. This pattern is consistent with a practical saturation point: a
small amount of causal temporal structure appears valuable, but the additional complexity of
R2e has not been shown to yield a reliable patient-level gain over R1a.

The unconstrained R3 regressors illustrate why parameter count alone is not the relevant
comparison. Their apparent internal advantages varied by reference, and they showed variance
shrinkage and weaker performance on the trigger-based evaluation. On the independent Stage~4
phantom, their in-band scenario RMSEs were 3.18--3.64~bpm, while R2e's was 0.0197~bpm with a
calibration slope of 0.9994. R3a and R3b used ten training infants, with one infant reserved for
checkpoint selection; R3c and R2e used all eleven. The training difference is a fairness
limitation, although R3c also showed substantial phantom error. The results therefore support
physical and shift-equivariant constraints over these tested shortcuts in these evaluations, not
a general conclusion that neural models cannot transfer. Shortcut learning is a recognized risk
when models can exploit correlations that do not persist under shift \citep{geirhos2020shortcut}.

The independent lines of evidence are complementary rather than interchangeable. Stage~2 tests
patient-held-out estimation on NICU recordings; Stage~4 tests frozen algorithms on controlled
stationary phantom motion. The Stage~4 benchmark comprised twenty conditions, fourteen in the
filter's in-band range and six below it. R2e performed exceptionally on the stable in-band
phantom signals, but Stage~4 contained no infants, spontaneous motion, clinical care, or bedside
reference uncertainty. Its result supports the project's L4 claim of zero-touch physical
transfer under the tested protocol. It cannot be used as evidence of L5 clinical utility.

\subsection{Uncertainty is useful only when its operating limits are visible}

The uncertainty analyses do not establish a second, independently validated clinical benefit.
U2, posterior standard deviation, is attractive because it is a direct posterior-spread score
with no fitted parameters or additional cutoff. Its session-level E-AURC was
$-0.2454$~bpm (95\% interval [$-0.4069,-0.1168$]) on the monitor reference and
$-0.2929$~bpm ([$-0.4555,-0.1436$]) on the trigger reference. At 80\% retained coverage it
reduced the one-hertz session-macro RMSE from 6.8486 to 6.6890~bpm; 10/11 infants improved. The raw
patient-level $p$-value was 0.0098, but the Holm-adjusted value was 0.0781. On the trigger
reference, 9/11 infants improved and the raw $p$-value of 0.0244 became 0.1953 after Holm
correction. The negative E-AURC estimates and these patient-level patterns are promising
internal evidence, not multiplicity-adjusted confirmation that U2 improves clinical decisions.

U6b addresses a different operating need: it combines posterior and spectral information to
flag questionable signal windows while limiting rejection during quiet pause-like intervals.
Its weights are not fitted, but it uses the fixed, empirically informed
$\tau_{\mathrm{ood}}=0.90$ cutoff and a three-hop (6~s) history. At that setting the pause-rejection
ratio was 1.1275 pooled and 1.1209 session-macro, within the project's 1.15 criterion; the
synthetic white-noise AUROC was 0.9986, with a 0.0036 false-positive rate at 95\% sensitivity.
On Stage~4, U6b separated the tested 5--10~bpm stopband conditions from in-band phantom motion
with AUROC 0.9998. When the 20~bpm skirt was added to the out-of-band class, its AUROC fell to
0.6686. Accordingly, the evidence supports a useful but fixed heuristic for the tested pause
and corruption conditions. It does not establish a generally reliable OOD detector or a
threshold-free method.

The learned five-parameter U6c score provides a counterexample to treating selective-risk
performance as the only criterion: although its risk--coverage point estimates were favorable,
its pause-rejection ratio of 1.2266 exceeded the project's 1.15 limit. Conversely, U6b's safer
pause ratio does not make it universally superior. Score ranking, pause preservation, corruption
detection, and patient-level uncertainty are separate properties and should remain separately
reported. In this cohort, a safe monitor design would expose quality or uncertainty to the
operator without silently suppressing a potentially important pause.

\section{The gap between phantom transfer and clinical readiness}
\label{sec:ch7_readiness}

The clearest quantitative readiness gap is the difference between Stage~4 point accuracy and
Stage~3 patient-level uncertainty. R2e's 0.0197~bpm scenario-macro RMSE on the in-band phantom
reflects a stable periodic signal with a laser reference. In the clinical cohort, the locally
adaptive C5 interval procedure at an 80\% target had a mean full width of 24.02~bpm. Its mean
patient coverage was 0.864, but the minimum was 0.586 and two of eleven infants were below the
nominal target. The full width averaged 17.96~bpm in the lower-posterior-uncertainty group and
30.54~bpm in the higher-uncertainty group. These widths are cohort-specific empirical results;
they are not half-widths, and they do not establish a universal physical limit of 24~GHz radar.
They show that, for this model, front end, cohort, and calibration design, intervals wide enough
to cover patient-level errors can be poorly suited to narrow rate decisions.

The finite-sample rank constraint is more fundamental to the current 95\% calibration claim. In
each leave-one-patient-out fold, only ten infants are available for calibration. The patient-level
order-statistic rank at 95\% is
\[
  \left\lceil (K_{\mathrm{cal}}+1)(1-0.05)\right\rceil
  =\lceil 11\times0.95\rceil=11,
\]
which exceeds the ten observed calibration units. A finite patient-exchangeable 95\% interval
therefore cannot be returned by this rank construction; its set is vacuous. At least
$K_{\mathrm{cal}}=19$ independent calibration patients are required for a finite 95\% interval
under this rule. That is a mathematical minimum for rank resolution, not a sufficient sample
size for clinical validation, subgroup precision, or safe alarm evaluation. Repeated windows
from the same infant do not increase the number of independent calibration patients. This
distinction between window-level abundance and patient-level information is central to
interpreting conformal intervals in clustered clinical data \citep{dunn2023distribution,barber2023conformal}.

These interval results also need careful interpretation at the 80\% target. C5 is a two-level,
locally adaptive calibration procedure informed by hierarchical conformal methods. If the same
nominal level is used at both levels, the simple hierarchical coverage bound under the required
assumptions is $(1-\alpha)^2=0.64$, not 0.80; the present analysis does not establish those
assumptions or a distribution-free, per-infant 80\% guarantee \citep{dunn2023distribution}.
Its reported coverage is empirical and heterogeneous across infants. Wider calibration sets or
a different cohort could change both coverage and width. The 24.02~bpm mean full width is
therefore best treated as a readiness diagnostic for the present single-channel pipeline, not
as a sensor specification or immutable information-theoretic bound.

The NICU recordings also contain plausible sources of ambiguity that are not represented by a
stationary phantom. At the session level, the fraction of windows in the cardiac HR/2 collision
band was associated with excess respiratory-rate error (Spearman $\rho=+0.696$,
$p<10^{-4}$). This association is consistent with harmonic ambiguity but does not prove that
cardiac motion caused the error. Session~11.1 had an I/Q standard-deviation ratio of 5.86,
absolute I/Q correlation of 0.863, and clipping in 5.06\% of samples; its shared-lag RMSE was
12.65~bpm. Session~9.2 also contained intervals where the recorded UCD status was not
\texttt{ChildPresent}. Random body movement and receiver or firmware interruptions add further
variation. These examples document observed failure contexts; they do not provide a causal
decomposition of the whole cohort.

The target reference itself is a limitation. The 1~Hz monitor stream and the trigger-derived
rates do not measure exactly the same quantity. One comparison excludes inter-trigger gaps over
five seconds, while a whole-window count includes the pause duration; these pause-like windows
are not clinically adjudicated apnea diagnoses. Differences between references, monitor delay,
and synchronization can all affect measured error. A prospective evaluation will need a
prespecified reference protocol, explicit timing alignment, and a clear treatment of pauses,
missingness, and alarm events. Neither the present internal references nor the stationary
phantom can show that a radar alarm would improve outcomes or safely replace a wired monitor.

Label-dependent online conformal adaptation does not remove this gap by itself. In standard
adaptive conformal inference, the update uses the realized miss indicator
$e_t=\mathbb{1}\{y_t\notin\widehat{C}_t(\alpha_t)\}$ to modify the next miscoverage level, for example
\[
  \alpha_{t+1}=\alpha_t+\gamma(\alpha-e_t).
\]
The indicator cannot be computed until a reference respiratory rate is observed. If a bedside
system has no concurrent reference label, this feedback update is unavailable; it cannot
autonomously adapt its interval from the radar stream alone. ACI and related online methods remain
possible in a supervised shadow study, or with delayed and independently measured labels, but
their label requirement and operational delay must be part of the design \citep{gibbs2024conformal,angelopoulos2024conformal}.

\section{Limitations and claim boundaries}
\label{sec:ch7_limitations}

Several limitations qualify the conclusions. First, the internal cohort contains eleven infants
and 28 synchronized sessions. The infant, rather than the one-hertz window, is the relevant
independent unit for claims about performance across patients. This sample constrains power for
close model comparisons, subgroup analysis, and patient-level conformal calibration. In
particular, failure to detect an R2e--R1a difference is not proof of equivalence.

Second, the references are imperfect and operationally distinct. The monitor trend has temporal
smoothing and delay; respiratory-trigger rates depend on trigger timing and on whether long gaps
are omitted from the rate denominator. The chapter results consequently answer a comparison
question under defined reference constructions, not an absolute measurement-accuracy question
against a universally valid gold standard. Patient-level analyses reduce the risk of treating
correlated windows as independent evidence, but they cannot resolve uncertainty in the reference
labels themselves.

Third, the acquisition and front-end choices bound the observable rate range. The locked
64-tap zero-phase Kaiser FIR passband is 0.30--3.00~Hz. The 5--10~bpm phantom conditions lie
below its lower edge, and rates in that region cannot be inferred reliably from the current
filtered representation. Raising the low-frequency response or changing the window and
separation stages could extend the range, but those changes would alter the frozen pipeline and
require new tests of latency, motion artifacts, harmonic confusion, and patient-level safety.
The present work does not validate severe bradypnea detection.

Fourth, some uncertainty design choices remain empirical. U6b's fixed cutoff of 0.90 and
three-hop history satisfy the selected pause and white-noise criteria on this cohort, while
nearby choices reveal a trade-off. The HQI harmonic-deficit normalization also compresses its
observed range and is highly correlated with an alternative non-overlapping denominator
(Spearman $\rho=0.999945$); this is a reason to interpret its scale cautiously, even though the
ranking is similar. These engineering decisions should be frozen before external clinical
testing and assessed across receiver settings and care conditions.

Fifth, Stage~4 is deliberately narrow. It tests stationary sinusoidal phantom movement on
independent 24.2~GHz hardware and supports a zero-touch L4 transfer result for the tested
protocol. The performance does not establish transfer to infants, across institutions, or to
different radar placements and waveforms. Exact reproduction of the original Beltr\~ao et al.
clinical subset also remains unresolved because cohort identifiers, the precise denominator,
and parts of the synchronization and solver protocol are unavailable. The work therefore does
not claim exact reproduction of that study (L2), nor clinical bedside utility, autonomous apnea
alarming, or superiority to wired NICU monitoring (L5) \citep{beltrao2022contactless}.

These boundaries are consistent with broader reviews of neonatal non-contact and wireless
monitoring, which report small studies, limited evaluation of feasibility and safety, and
methodological gaps alongside technical promise \citep{senechal2026next}. Such context does not
validate or invalidate the present estimator; it reinforces the need to evaluate signal
availability, safety, reference quality, and workflow in addition to point accuracy.

\section{Next research steps}
\label{sec:ch7_future}

The next phase should address both sensing ambiguity and the lack of independent patient-level
evidence. On the engineering side, multi-antenna or wideband FMCW measurements could provide
range and spatial information that a single-channel CW system does not observe. Range-bin
isolation may help distinguish thoracic motion from other moving reflectors, while optical or
time-of-flight channels could provide complementary motion information. These are research
directions, not guaranteed fixes: additional channels introduce calibration, synchronization,
privacy, and device-integration requirements. Recent neonatal radar and multimodal studies
provide relevant starting points for controlled comparisons \citep{liao2025noncontact,hari2025contactless,gleichauf2021automated}.

The statistical work should first expand the number of independent infants. Nineteen calibration
patients is the rank-based minimum for a finite patient-level 95\% interval under the stated
order-statistic rule; a larger multicenter cohort will be needed to estimate performance across
gestational ages, body sizes, clinical support, motion states, and sites. The final sample size
should follow a prespecified precision or power calculation and account for unusable recordings
and site-level variation. Calibration and evaluation must be separated at the patient and,
where appropriate, site level. Any proposed conformal procedure should be audited for its actual
exchangeability unit, finite-sample rank, coverage by patient and subgroup, interval width, and
behavior under distribution shift rather than relying on a stored validity flag alone.

Clinical testing should proceed prospectively in shadow mode while existing monitoring remains
available. The protocol should define the intended use before data collection: respiratory-rate
trend display, quality warning, or candidate alarm support are different functions and require
different endpoints. It should measure usable-data availability, time alignment, error and
coverage by infant, false and missed alarms, alert delay, pause handling, and the burden of
non-actionable alerts. Any transition toward interventional use would require ethics review,
appropriate parent or guardian consent, privacy and data-governance protections, and a clear
clinical oversight and escalation pathway. The safety case must make abstention visible and must
not turn a rejected radar window into silence when the underlying event could be clinically
important.

\section{Ethics, reproducibility, and conclusion}
\label{sec:ch7_ethics_conclusion}

The principal ethical constraint is to avoid presenting a research estimate as a replacement for
clinical observation. In particular, a quality filter that rejects a pause-like window can hide
the very period that requires review. Future interfaces should distinguish a low-confidence
estimate from a low respiratory rate, preserve access to the underlying signal and quality
status, and defer clinical action to an approved workflow. Data use, access, storage, and any
future model update must follow the applicable institutional and clinical governance. The use
of Paper Orchestra and other AI-assisted drafting or editorial tools should be disclosed
in the required authorship declaration, including the tools and scope of assistance and the
distinction between editorial support and scientific analysis. Responsibility for the scientific
claims, analyses, and final text remains with the author.

The project records a frozen analysis basis at commit \texttt{5e458c1}, with the legacy baseline
environment identified as Python~3.9.6 and the WP1--WP5 research pipeline as Python~3.10.19.
The reproducibility appendix reports integrity checks for the archived scripts, caches,
figures, bibliography, and notation contract. Those records support review of how the reported
results were produced; they do not substitute for independent replication or release of data
that cannot ethically or legally be shared.

In conclusion, this thesis finds that a compact, physics-constrained causal estimator can
outperform a memoryless baseline across two defined NICU references and can transfer with very
low error to an independent stationary phantom. It also shows that posterior uncertainty and a
carefully selected quality barrier can help characterize when estimates may be unreliable, while
finite patient counts impose sharp limits on interval resolution. The evidence is strongest for
the value of explicit structure and disciplined evaluation within the tested system. It remains
insufficient for a claim of prospective bedside utility. The responsible next step is larger,
patient-level, multicenter validation with transparent uncertainty, preserved standard care, and
predefined safety outcomes.
